%% LyX 2.4.4 created this file.  For more info, see https://www.lyx.org/.
%% Do not edit unless you really know what you are doing.
\documentclass[english]{article}
\usepackage[T1]{fontenc}
\usepackage[latin9]{inputenc}
\usepackage{amsmath}
\usepackage{amsthm}
\usepackage{amssymb}
\usepackage{geometry}
\geometry{verbose,tmargin=1in,bmargin=1in,lmargin=1in,rmargin=1in}

\makeatletter
%%%%%%%%%%%%%%%%%%%%%%%%%%%%%% User specified LaTeX commands.
\usepackage{fancyhdr}
\usepackage{lastpage}
\pagestyle{fancy}
\usepackage{enumitem}
\usepackage{tikz}
\usetikzlibrary{shapes,snakes}
\usepackage{amsmath,amssymb}
\usepackage{multicol}
% Added by lyx2lyx
\setlength{\parskip}{\smallskipamount}
\setlength{\parindent}{0pt}

\makeatother

\usepackage{babel}
\begin{document}
\lhead{MTH 331(A) Dr.\,James Ashe}
\chead{\textbf{Helpful Results}}
\cfoot{\thepage{}~of~\pageref{LastPage}}
%%
%%\setenumerate[0]{leftmargin=12pt,itemindent=0pt,labelwidth=0pt}
\renewcommand{\labelenumi}{\textbf{\arabic{enumi}.}}

\subsection*{Chapter 9}

\textbf{Squeeze Theorem for Sequences. }Let $\{a_{n}\}$, $\{b_{n}\}$,
and $\{c_{n}\}$ with $a_{n}\leq b_{n}\leq c_{n}$ for all integers
\emph{n }greater than some index $N$. If ${\displaystyle \lim_{n\rightarrow\infty}a_{n}}={\displaystyle \lim_{n\rightarrow\infty}b_{n}=L}$,
then $\lim_{n\rightarrow\infty}c_{n}=L$.\vfill

\textbf{Growth Rate of Sequences. }The following sequences are ordered
according to increasing growth rates as $n\rightarrow\infty$: $\{\ln n\}\ll\{n^{p}\}\ll\{n^{p}\ln n\}\ll\{n^{p+s}\}\ll\{b^{n}\}\ll\{n!\}\ll\{n^{n}\}$.
The ordering applies for positive real numbers $p,r$, and $s$ and
$b>1$.\vfill

\textbf{Geometric Series. }Let $r,a\in\mathbb{R}$. If $|r|<1$, then
${\displaystyle \sum_{k=0}^{\infty}}ar^{k}=\frac{a}{1-r}$. If $|r|\geq1$,
then the series diverges. \vfill

\textbf{Divergence Test. }If $\sum a_{k}$ converges, then ${\displaystyle \lim_{n\rightarrow\infty}a_{k}=0}.$
Equivalently, if ${\displaystyle \lim_{n\rightarrow\infty}a_{k}\neq0}$,
then the series diverges. \vfill

\textbf{Integral Test. }Suppose $f$ is a continuous, positive, decreasing
function for $x\geq1$ and let $a_{k}=f(k)$ for $k=1,2,3,\dots$.
Then 
\[
{\displaystyle \sum_{k=1}^{\infty}a_{k}}\mbox{ and \ensuremath{\int_{1}^{\infty}f(x)\,dx}}
\]
either both converge or both diverge. In the case of convergence,
the value of the integral is \emph{not}, in general, equal to the
value of the series. \vfill

\subsection*{Chapter 10}

\textbf{Taylor's Polynomials. }Let $f$ be a function with $f',f'',\dots f^{(n)}$
defined at $a$. The \textbf{\emph{n}}\textbf{th-order Taylor Polynomial
}for $f$ with its \textbf{center }at $a$ denoted $p_{n}$, has the
property that it matches $f$ in value, slope, and all derivatives
up to the \emph{n}th derivative at $a$; that is,
\[
p_{n}(a)=f(a)\mbox{ and }p'_{n}(a)=f'(a),\dots,p_{n}^{(n)}(a)=f^{(n)}(a)\mbox{.}
\]

The \emph{n}th-order Taylor polynomial centered at $a$ is 
\[
p_{n}(x)=f(a)+f'(a)(x-a)+\frac{f''(a)}{2!}(x-a)^{2}+\dots+\frac{f^{(n)}(a)}{n!}(x-a)^{n}\mbox{.}
\]

More compactly, 
\[
p_{n}(x)={\displaystyle \sum_{k=0}^{n}\frac{f^{(k)}(a)}{k!}(x-a)^{k}\mbox{ for }k=0,1,2,\dots,n\mbox{.}}
\]
\vfill

\subsection*{Helpful formulas}

\begin{multicols}{2}
\begin{enumerate}
\item $\frac{d}{dx}\left(x^{n}\right)=nx^{n-1}$
\item $\frac{d}{dx}\left[f\left(g(x)\right)\right]=f'\left(g(x)\right)\cdot g'(x)$
(Chain rule)
\item $\frac{d}{dx}\left(f(x)g(x)\right)=f'(x)g(x)+f(x)g'(x)$
\item $\frac{d}{dx}e^{x}=e^{x}$
\item $\frac{d}{dx}\sin ax=a\cos ax$
\item $\frac{d}{dx}\cos ax=-a\sin ax$
\item $\int x^{n}\,dx=\frac{1}{n+1}x^{n+1}+C$; $n\neq1$ 
\item $\int\cos ax\,dx=\frac{1}{a}\sin ax+C$
\item $\int\sin ax\,dx=-\frac{1}{a}\cos ax+C$
\item $\int\cos^{2}x\,dx=\frac{x}{2}+\frac{\sin2x}{4}+C$
\item $\int\sin^{2}x\,dx=\frac{x}{2}-\frac{\sin2x}{4}+C$
\end{enumerate}
\end{multicols}
\end{document}
